\section{The large relative gaps}
\label{app:positive-boundary}

We separate the prime-power gaps from the three contour estimates.
The arithmetic proof for prime-power gaps is given below. The
electronic appendix \textit{Positive-exponent Boundary} gives the
quantitative contour bounds, their derivations and the exact
constant checks.

\begin{lemma}
\label{lem:positive-boundary}
The pair \(P_{N-1,a,d},P_{N,a,d}\) is coprime for every \(N\ge2\)
on \(a=1,d\ge3\), and on \(a\ge2,d\ge3a\).
\end{lemma}

\begin{proof}
Proposition~\ref{prop:small-terminals} covers \(N\le24\).
Thus we may impose the ratio bound of Lemma~\ref{lem:ratio-range}
throughout this proof.

For \(a=1\), write
\[
 Q_{m,\gamma,d}(x)=\sum_{k=0}^m\binom{m}{k}
                        (\gamma)_{dk}x^{m-k}.
\]
At terminal \(N\) the actual shapes are \(N\) and \(N+1\).
Coefficient comparison gives, with \(\gamma=N\),
\[
 Q_{N,\gamma,d}'=NQ_{N-1,\gamma,d},\qquad
 \gamma Q_{N,\gamma+1,d}
  =(\gamma+dN)Q_{N,\gamma,d}-dxQ_{N,\gamma,d}'.
\]
In fact the lower reciprocal is \(Q_{N-1,N,d}\), so at a
common root the derivative of \(Q_{N,N,d}\) vanishes. The second
identity then makes the upper reciprocal a nonzero scalar multiple
of \(Q_{N,N,d}\) there. Conversely, a multiple root of
\(Q_{N,N,d}\) makes both reciprocals vanish. The prime-power
argument in the next subsection gives squarefreeness at
\(d=3,4,5,7,8,9\). The separate
gap-six contour proof gives it at \(d=6\) for \(N\ge10000\).

For \(d\ge10\), the lower-one contour proof gives squarefreeness of
\(Q_{N,N,d}\) whenever \(N\ge\max(5000,128(d+1))\).
Its two contour contributions have separated logarithmic
derivatives at a zero, so the derivative of their sum is nonzero.
For \(a\ge2,d\ge3a\), the weighted contour proof in the electronic
appendix gives coprimality for \(N\ge\max(4000,128(a+d))\).
Here one must compare two integrals with the actual moment orders.
With \(w=A/(BN^d)\) and \(g(x)=x^a(w+x^d)\), their integrands
after the common dilation are \(g(x)^N e^{-Nx}\) and
\(g(x)^{N-1}e^{-Nx}\). Write \(I_U,I_L\) and \(J_U,J_L\)
for their two contour contributions. At a zero of the first
integral, \(I_U=-I_L\ne0\). The verified packet estimates at the
two selected saddles \(x_U,x_L\) give
\[
 \left|\frac{J_U}{I_U}\right|
 <\frac{24\sqrt{N/(N-1)}}{|g(x_U)|},\qquad
 \left|\frac{J_L}{I_L}\right|
 >\frac{\sqrt{N/(N-1)}}{5|g(x_L)|},\qquad
 \left|\frac{g(x_U)}{g(x_L)}\right|>150 .
\]
Thus \(|J_U|<|J_L|\) and the second integral is nonzero.
The proof estimates the extra amplitude \(e^{-x}\) in the
order-\(N-1\) contour integral when the phase is written as
\(\log g(x)-x\).

If \(N<\max(4000,128b)\), where \(b=a+d\le4d/3\), a
potential failure has \(d\le24\). Indeed \(s\ge3d/4\).
For \(N\ge200\) and \(d\ge25\), one has
\[
 N<512d/3,\qquad \frac{3d}{4}>2\log(512d/3)+2 .
\]
At \(d=25\) the margin exceeds \(3/100\), and its derivative
\(3/4-2/d\) is positive. For \(N<200\), use
\(\sigma_N\le 2\log 199+3<14<3d/4\).

For \(a=1,d\ge10\) below the contour threshold,
a potential failure has \(d\le19\).
For \(20\le d\le38\) the threshold is \(5000\), and
\[
 \frac{d^2}{d+1}\ge\frac{400}{21}>2\log(5000)+2 .
\]
For \(d\ge39\) the threshold is \(128(d+1)\). The difference
\(d^2/(d+1)-2\log(128(d+1))-2\) exceeds \(18\) at \(39\),
and its derivative \(1-1/(d+1)^2-2/(d+1)\) is positive.
For \(N<200\) and \(d\ge20\), the upper bound
\(\sigma_N<14<400/21\) suffices. The electronic comparison
script certifies the displayed logarithmic inequalities by
rational bounds.

The supplemental finite checker therefore needs only the families
\[
\begin{array}{c|c|c}
a&d&\text{terminals }N\\ \hline
1&6&25\le N<10000\\
1&10\le d\le19&25\le N<5000\\
2\le a\le\lfloor d/3\rfloor&6\le d\le24&
                25\le N<\max(4000,128(a+d)).
\end{array}
\]
Before any further exclusions this is \(338\,071\) cases.
Every case satisfies \(a<d(N-2)\), and a possible common factor
has \(s<21\) by Lemma~\ref{lem:ratio-range}. The block-prime
support and Newton allocation arguments in
Lemmas~\ref{lem:prime-support} and~\ref{lem:newton-alloc}
apply here with these hypotheses. The stronger degree bound
is used only when its starting-point condition is verified
for that degree. In the subsequent Newton check, each odd
prime \(p>aN\) satisfies \(p>N,p>d\) and \(p^2>(a+d)N\)
throughout these finite families. Thus the sparse floor
ordinates used in Section~\ref{sec:arithmetic} are exact
also in this range.

The supplemental program \texttt{verify\_positive\_boundary.cpp}
retains \(326\,146\) cases after imposing the uniform bound
\(d^2<21(a+d)\). It tries the
first available sparse Newton test, applies the prime-block
bounds at the remaining degrees, and then uses additional
primes as needed. The exact exclusions are
\[
\begin{array}{l|r|r|r}
\text{Family}&\text{Supports}&\text{Prime blocks}&\text{Newton polygons}\\ \hline
a\ge2,\ d\ge3a &266\,421&12\,901&253\,520\\
a=1,\ 10\le d\le19&49\,750&966&48\,784\\
a=1,\ d=6&9\,975&535&9\,440
\end{array}
\]
Every possible even degree is excluded, and no modular fallback
is needed.
\end{proof}

\subsection{The prime-power gaps}

We prove the squarefreeness needed above for \(q=p^r\),
where \(p\) is prime and \(r\ge1\). Use \(n=N\), \(\gamma=n\)
and
\[
 T(t)=t^n Q_{n,\gamma,q}(1/t),\qquad
 R(z)=z^{qn}T((-1)^{q+1}z^{-q}),\qquad M=\gamma+qn.
\]
Both changes of variable are unramified away from zero, and both
\(T(0)\) and \(R(0)\) are nonzero. Write \(D=d/dz\), \(\theta=zD\)
and define
\[
 \mathcal A=z-zD+\gamma+q-1,\qquad
 \Pi_{n,q}(t)=\prod_{j=0}^{n-1}(t+qj),\qquad
 B_n=\Pi_{n,q}(\mathcal A)1 .
\]

\paragraph{The rank condition.}
The polynomial \(R\) solves \(LR=0\) for the monic operator
\[
 L=z^{-q-1}\left[
 \theta\prod_{\ell=0}^{q-1}(\theta-M-\ell)-z^q(\theta-qn)\right].
\]
Indeed consecutive coefficients \(r_k=[z^{q(n-k)}]R\) satisfy
\((k+1)r_{k+1}=(-1)^{q+1}(n-k)(\gamma+qk)_q r_k\),
which is the equation at \(z^{q(n-k)}\) after multiplying by
\(z^{q+1}\). Put \(U=1-D\). We will use the adjoint solutions
\[
 f_\omega(z)=e^{-\omega z}z^{1-M}B_n(\omega z),
 \qquad \omega^q=1.
\]
To derive them, set
\[
 \begin{aligned}
 C(x)&=(1-x)^{-\gamma}
       \left(\frac{1-(1-x)^q}{qx}\right)^{-n-1},\\
 A_0(x)&=(1-x)((1-x)^q-1),\\
 b_0(x)&=(M-1)((1-x)^q-1)+qn.
 \end{aligned}
\]
The residue \(n![x^n](C(x)e^{zx})\) equals \(B_n\).
In fact the substitution \(v=1-(1-x)^q\) makes it
\[
 q^n n![v^n](1-v)^{-(\gamma+q-1)/q}
                  e^{z(1-(1-v)^{1/q})}.
\]
With \(v=qt\) the generating function in \(n\) is
\((1-qt)^{-(\gamma+q-1)/q}e^{z(1-(1-qt)^{1/q})}\).
Writing it as \(F(t,z)\), differentiation gives
\((1-qt)F_t=\mathcal A F\), \(F(0,z)=1\). This proves the
claimed residue formula. Logarithmic differentiation gives
\(A_0 C'=((n+1)A_0/x-A_0'+b_0)C\). Taking the residue of
\((e^{zx}A_0 C x^{-n-1})'\) therefore yields
\[
 [zA_0(D)+b_0(D)]B_n
 =z(U^{q+1}-U)B_n+(M-1)(U^q-1)B_n+qnB_n=0.
\]
The divergence form \(L=z^{-q}D z^{M+q}D^q z^{-M}-D+qn/z\)
gives \(L^*=(-1)^{q+1}z^{-M}D^q z^{M+q}D z^{-q}+D+qn/z\).
For \(f_1=e^{-z}z^{1-M}B_n\) the first term becomes
\(e^{-z}z^{-M}[zU^{q+1}+(M-1)U^q]B_n\), using
\(U^qz=zU^q-qU^{q-1}\). The last two terms give
\(e^{-z}z^{-M}[-zU+qn+1-M]B_n\).
This proves \(L^*f_1=0\), and substitution of \(\omega z\)
gives the remaining solutions.

Write \(L=\sum_{k=0}^{q+1}a_kD^k\). The bilinear expression
\[
 \sum_{k=1}^{q+1}\sum_{i=0}^{k-1}
 (-1)^i(a_k f_\omega)^{(i)} R^{(k-1-i)}
\]
has derivative \(f_\omega LR-RL^*f_\omega=0\), by cancellation of
adjacent summands. It is zero because it is a rational function of
\(z\) times \(e^{-\omega z}\), which cannot be a nonzero constant.
At a multiple root \(z_0\ne0\) of \(R\), the first two coordinates
of its jet vanish, and its coordinates \(R'',\ldots,R^{(q)}\)
are not all zero by uniqueness for the ordinary equation \(LR=0\).
They give a common nonzero nullvector of the restricted rows of
these expressions. In the column multiplying \(R^{(j)}\), \(2\le
j\le q\), the term of highest derivative of \(f_\omega\) is
\((-1)^{q-j}f_\omega^{(q-j)}\), since \(a_{q+1}=1\).
The restricted rows are therefore triangular transforms of the
jets \(f_\omega,\ldots,f_\omega^{(q-2)}\). Removing the common
power gauge and the nonzero exponential row factors gives a column
dependence among \(\omega^j(U^jB_n)(\omega z_0)\), \(0\le j\le q-2\).
The relation \(Uz=zU-1\) shows that \((zU)^j\) differs from
\(z^jU^j\) by a combination of \(z^iU^i\), \(i<j\).
Using \(\mathcal A=zU+\gamma+q-1\) and \(z_0\ne0\), the dependence
becomes
\[
 [\Pi_{n,q}(\mathcal A)P(\mathcal A)1](\omega z_0)=0
 \quad\text{for every }\omega^q=1
\]
for some nonzero \(P\) of degree at most \(q-2\).
The dependent evaluation matrix has algebraic entries, so this
condition transfers to \(\overline{\Q}_p\).

\paragraph{Excluding the dependence.}
Normalize \(v(p)=1\), work in a complete algebraically closed
extension of \(\Q_p\), and let \(\mathcal O\) be its valuation ring.
Set \(b_r=1+p+\cdots+p^{r-1}\). Every interval of \(qk\) consecutive
integers contains \(qk/p^\ell\) multiples of \(p^\ell\) for
\(1\le\ell\le r\). Thus \(v((\gamma)_{qk})\ge b_r k\), and
\(v(z_0)\ge b_r/q\). Otherwise the leading term of \(R(z_0)\)
would have uniquely least valuation. In particular
\(v(u)\ge b_r\) for \(u=z_0^q\).

Here is the invariant under descent from \(q=p^r\) to \(p^{r-1}\).
Allow an operator and seed
\[
 A=c-zD+z+\sum_{i\ge1}a_i z^{p^i},\quad
 c\in\Z_p,\quad a_i\in p\Z_p,\qquad
 h\in\mathcal O[[z]],\quad \bar h=\overline{h(0)}\ne0 .
\]
For \(u\ne0\), \(v(u)\ge b_r\), and primitive \(P\in\mathcal O[t]\)
with \(\deg\bar P\le q-2\), the cyclic coefficient vector of
\(\Pi_{n,q}(A)P(A)h\) modulo \(z^q-u\) is nonzero.
Here the component \(j\) of the vector of a series \(\sum f_kz^k\)
is \(\sum_{k\ge0}f_{qk+j}u^k\). These sums converge since \(v(u)>0\),
and their simultaneous vanishing is equivalent to vanishing at
the \(q\) distinct roots of \(z^q=u\).

At \(r=1\) reduce the vector modulo the maximal ideal.
The reduced ring is \(k[z]/(z^p)\), where \(k\) is the residue field,
\(h\) reduces to a constant unit, and \(\Pi_{n,p}(\bar A)=\bar A^n\).
The matrix of \(\bar A=z-zD+\bar c\) is triangular with distinct
diagonal entries \(\bar c-j\), \(0\le j<p\). They are all the
elements of \(\mathbb F_p\). Its vector \(1\) is cyclic because
\(\bar A^j1\) is monic of degree \(j\) for \(j<p\), so \(1\) has
a nonzero projection to each eigenspace. Vanishing would force
\(\bar P\) to vanish at all \(p-1\) nonzero eigenvalues, impossible
by its degree.

At \(r\ge2\) use the integral unit series
\[
 E(z)=\exp\!\left(\sum_{i\ge0}z^{p^i}/p^i\right)
     =\prod_{p\nmid m}(1-z^m)^{-\mu(m)/m},
\]
where \(\mu\) is the M\"obius function. The product identity follows
by taking formal logarithms and summing \(\mu\) over divisors.
Each factor and its inverse are integral because their exponents
are in \(\Z_p\) and the binomial coefficients are integral there
by continuity from nonnegative integers. The conjugation \(E^{-1}AE\)
subtracts \(zD\log E=\sum_{i\ge0}z^{p^i}\) from \(A\).
It preserves vanishing after changing the seed to \(E^{-1}h\),
since the series and their inverses converge at every point of the
fiber. In residue \(j\bmod p\), set \(y=z^p\),
\(b=a_1-1\), and \(w=by/p\). Write \(e_j\in\{0,\ldots,p-1\}\)
for the representative of \(c-j\bmod p\). The new block operator
is \(e_j+pL_j\), with
\[
 L_j=\frac{c-j-e_j}{p}-w\frac d{dw}+w+
        \sum_{i\ge1}\frac{(a_{i+1}-1)p^{p^i-1}}{b^{p^i}}w^{p^i}.
\]
This is in the same class. If \(E^{-1}h=\sum_{j=0}^{p-1}z^j
\widetilde h_j(y)\), the new seed
\(h_j(w)=\widetilde h_j(pw/b)\) reduces to the constant
\(\overline{h(0)}(-1)^j/j!\ne0\), using the first \(p\)
coefficients of \(e^{-z}\). Put \(q'=q/p\). The new fiber parameter
is \(u_*=(b/p)^{q'}u\), with \(v(u_*)\ge b_r-q'=b_{r-1}\).
If the original cyclic vector vanishes, the block vectors
\[
 \Pi_{n,q}(e_j+pL_j)P(e_j+pL_j)h_j
             \pmod{w^{q'}-u_*}
\]
all vanish. Indeed for each \(y^{{q'}}=u\), the evaluations on the
\(p\) roots of \(z^p=y\) separate these \(p\) residue blocks by a
Vandermonde matrix.

If \(e_j\ne0\), absorb the first factor into the seed \(\widehat h_j\).
This still reduces to a constant unit. The cyclic vectors of
\(L_j^i\widehat h_j\), \(0\le i<q'\), form a unit basis since
their reductions are polynomials of successively increasing degree
with unit leading coefficient in \(k[w]/(w^{q'})\).
The higher vectors are integral. Expanding \(P(e_j+pL_j)\)
in divided Taylor coefficients \(P^{[i]}(e_j)\) and inverting
the unit basis gives \(v(P^{[i]}(e_j))\ge q'-i>0\) for \(i<q'\):
all terms with \(i\ge q'\) are divisible by \(p^{q'}\).
Hence
\[
 \bar P(t)=(t^{p-1}-1)^{q'}H(t),\qquad \deg H\le q'-2.
\]
The first unit coefficient of \(P\) occurs at an index at most
\(q'-2\). Choose a coefficient \(\kappa\) of \(P(pt)\) of
least valuation, so \(v(\kappa)\le q'-2\).
Then \(Q(t)=P(pt)/\kappa\) is primitive and \(\deg\bar Q\le q'-2\).
In the unique block \(e_j=0\),
\(\Pi_{n,q}(pL_j)=p^n\Pi_{n,q'}(L_j)\), and division by
\(p^n\kappa\) contradicts the induction hypothesis at \(q'\).
This proves the invariant.

Finally the original \(\mathcal A,1\) belong to this class.
Scaling the dependent polynomial \(P\) to be primitive, the
rank condition contradicts the invariant. This proves
squarefreeness at every prime-power gap needed in this appendix.
